\section{失败模式与工程防线}

\subsection{典型失败：误动作后无法恢复}
讲者在问答环节举了很典型的例子：模型执行了错误动作后，若缺少搜索和回溯，后续很难自我纠正，最终完全偏离任务目标。

\begin{figure}[H]
\centering
\includegraphics[width=0.82\textwidth]{frame-09-failure-mode.jpg}
\caption{视频约 01:17:20：讲者讨论错误动作后的恢复困难，强调搜索与回溯机制必要性}
\end{figure}

\begin{importantbox}{故障恢复的最小闭环}
\begin{enumerate}
    \item 检测：快速识别当前轨迹偏航；
    \item 回滚：回到最近可靠状态；
    \item 重规划：更新子目标与动作约束；
    \item 复核：通过 verifier 判定是否重回正轨。
\end{enumerate}
\end{importantbox}

\subsection{失败模式分类}
\begin{table}[H]
\centering
\small
\begin{tabular}{p{2.7cm}p{3.8cm}p{4.2cm}}
\toprule
\textbf{失败类型} & \textbf{典型症状} & \textbf{工程修复手段} \\
\midrule
感知错误 & 读错按钮/表单字段/图像内容 & 多模态校验、元素置信度阈值 \\
规划错误 & 子目标顺序错误或约束遗漏 & 计划模板、子目标检查器 \\
执行错误 & 连续误点、误输入、误停止 & 回放日志、状态快照与回滚 \\
评估错误 & Judge 误判成功/失败 & 多 Judge 交叉、人工抽检校准 \\
\bottomrule
\end{tabular}
\caption{Web Agent 常见失败模式与修复策略}
\end{table}

\begin{warningbox}{不要把 ``失败'' 当成噪声样本}
失败轨迹是最有价值的数据来源之一。系统应系统化归档失败类型，把它们喂回训练和评估流程，形成持续改进回路。
\end{warningbox}

\subsection{人类协同仍然必要}
对于高风险任务，Agent 应具备 ``请求澄清'' 和 ``升级到人工'' 的能力。人类不是最后兜底才出现，而应在不确定性升高时及时介入。

\begin{knowledgebox}{人机协同触发条件示例}
\begin{itemize}
    \item 连续两次高置信度分歧（planner 与 verifier 不一致）；
    \item 关键字段不可判定（支付、身份、隐私相关）；
    \item 预算超限但任务未收敛。
\end{itemize}
\end{knowledgebox}

\subsection{本章小结}
Agent 失败模式可被系统化管理。真正可靠的系统不是 ``永不出错''，而是 ``出错可检测、可恢复、可学习''。

